\section{拉萨危机与高原道路}

1705年拉藏汗推翻第巴桑结嘉措的政权安排，拉萨的政教权力结构发生剧变。《清圣祖实录》所见的是清廷收到的西藏信息；《西藏王臣记》、驻藏来文和旅行者材料则保存不同的时间线与行动者。清廷是否知情、支持到何种程度，不能从后来的军事介入倒推。第六世达赖喇嘛仓央嘉措于1706年被押离拉萨，途中死亡还是失踪仍有相互冲突的传统；这一不确定性不是叙事的空白，而是继承合法性本身被各方争夺的证据条件。

1708年，卡尔藏嘉措在青海塔尔寺附近被认定为转世之一（EQ-1708-KELZANG）。出生和早期认定大体可靠，具体仪式、时间与承认范围不能以晚出的转世谱系倒灌。青海寺院、蒙古关系和清廷册封语言共同构成了一个跨区域的合法性网络。1717年准噶尔军进入拉萨、杀拉藏汗（EQ-1717-DZUNGAR-LHASA），暴力规模、地方态度和准噶尔治理应并读藏文、满文与旅行记录，而不能由任何一方的“平定”或“解救”词汇覆盖。

清军1718年自四川方向入藏，在喀喇乌苏一带覆败（EQ-1718-QING-TIBET-EXPEDITION）。这一失败说明高原道路、气候、兵站和信息不足不是战报中的附注，而是跨区域权力的条件。满文奏折和藏文地名、地方记忆需彼此核对，特别不能把清方报告中的向导与伤亡解释为无争议事实。1720年清军与青海蒙古盟军进入拉萨、驱逐准噶尔力量并迎立卡尔藏嘉措为第七世达赖喇嘛（EQ-1720-LHASA-RESTORATION）。进入拉萨和迎立的主干可靠；盟军、西藏地方行动者与清军的角色不能被凯旋叙事抹平。军事存在的增强也不等于西藏政教制度被一道中央命令取代，后续治理仍须在多方关系中重新安排。


\subsection{事件链与材料限度}
青海与西藏的政教、道路和驻防关系需要多语材料相互限制。

\paragraph{康熙四十四年（1705年）：拉藏汗推翻第巴桑结嘉措的政权安排，西藏政教权力结构剧烈变化}
\eventanchor{EQ-1705-LHAZANG-TAKEOVER}{拉藏汗推翻第巴桑结嘉措的政权安排，西藏政教权力结构剧烈变化}。本节点叙述该事件本身。和硕特汗国、西藏噶厦与清在边地道路、驻防与多方社群交错的尺度的处境首先受第巴与和硕特汗权长期紧张，达赖喇嘛继承和对外联络加深拉萨政治危机制约。《清圣祖实录》康熙四十四年西藏条；《清史稿·西藏传》；《西藏王臣记》藏文与清廷驻藏来文留下可核的线索；可辨的后果是西藏内部权力重组为准噶尔介入和清廷后续军事行动创造条件，清方命令不等于拉萨实际控制。原有置信注记：政变主干可靠；清廷知情、支持程度和地方反应须并读藏文、满文与旅行者材料

\paragraph{康熙四十四年（1705年）：拉藏汗推翻第巴桑结嘉措的政权安排，西藏政教权力结构剧烈变化}
\eventanchor{EQ-1705-LHAZANG-TAKEOVER-AFTERMATH}{拉藏汗推翻第巴桑结嘉措的政权安排，西藏政教权力结构剧烈变化} 置于边地道路、驻防与多方社群交错的尺度中阅读，本节点追踪「拉藏汗推翻第巴桑结嘉措的政权安排，西藏政教权力结构剧烈变化」之后可见的余波。第巴与和硕特汗权长期紧张，达赖喇嘛继承和对外联络加深拉萨政治危机构成行动者和硕特汗国、西藏噶厦与清必须面对的条件。取证从《清圣祖实录》康熙四十四年西藏条；《清史稿·西藏传》；《西藏王臣记》藏文与清廷驻藏来文出发，因而只能把变化谨慎地连到西藏内部权力重组为准噶尔介入和清廷后续军事行动创造条件，清方命令不等于拉萨实际控制。原有置信注记：政变主干可靠；清廷知情、支持程度和地方反应须并读藏文、满文与旅行者材料

\paragraph{康熙四十四年（1705年）：拉藏汗推翻第巴桑结嘉措的政权安排，西藏政教权力结构剧烈变化}
\eventanchor{EQ-1705-LHAZANG-TAKEOVER-DECISION}{拉藏汗推翻第巴桑结嘉措的政权安排，西藏政教权力结构剧烈变化} 标记本节点定位围绕「拉藏汗推翻第巴桑结嘉措的政权安排，西藏政教权力结构剧烈变化」形成的处置与调度记录，涉及和硕特汗国、西藏噶厦与清。边地道路、驻防与多方社群交错的尺度不是抽象背景：第巴与和硕特汗权长期紧张，达赖喇嘛继承和对外联络加深拉萨政治危机。《清圣祖实录》康熙四十四年西藏条；《清史稿·西藏传》；《西藏王臣记》藏文与清廷驻藏来文所见的顺序随后指向西藏内部权力重组为准噶尔介入和清廷后续军事行动创造条件，清方命令不等于拉萨实际控制，但原有置信注记：政变主干可靠；清廷知情、支持程度和地方反应须并读藏文、满文与旅行者材料

\paragraph{康熙四十四年（1705年）：拉藏汗推翻第巴桑结嘉措的政权安排，西藏政教权力结构剧烈变化}
在边地道路、驻防与多方社群交错的尺度，\eventanchor{EQ-1705-LHAZANG-TAKEOVER-PRELUDE}{拉藏汗推翻第巴桑结嘉措的政权安排，西藏政教权力结构剧烈变化} 让和硕特汗国、西藏噶厦与清的一个选择或压力有了可查的坐标。本节点回看「拉藏汗推翻第巴桑结嘉措的政权安排，西藏政教权力结构剧烈变化」之前已可见的条件。前因可写为第巴与和硕特汗权长期紧张，达赖喇嘛继承和对外联络加深拉萨政治危机；《清圣祖实录》康熙四十四年西藏条；《清史稿·西藏传》；《西藏王臣记》藏文与清廷驻藏来文限定了记录，后续变化是西藏内部权力重组为准噶尔介入和清廷后续军事行动创造条件，清方命令不等于拉萨实际控制。原有置信注记：政变主干可靠；清廷知情、支持程度和地方反应须并读藏文、满文与旅行者材料

\paragraph{康熙四十五年（1706年）：第六世达赖喇嘛仓央嘉措被押离拉萨，途中死亡或失踪，其结局与继承合法性成为争议}
\eventanchor{EQ-1706-SIXTH-DALAI}{第六世达赖喇嘛仓央嘉措被押离拉萨，途中死亡或失踪，其结局与继承合法性成为争议} 所关联的是西藏、和硕特汗国与清在边地道路、驻防与多方社群交错的尺度中的行动。本节点叙述该事件本身。它从拉藏汗以宗教合法性重组拉萨权力，清廷册封语言与西藏各派、蒙古势力认定不一致展开，借《清圣祖实录》康熙四十五年西藏条；《清史稿·西藏传》；第巴传记、藏文史书与拉达克、蒙古相关记载进入可检验的年序，并以达赖喇嘛地位争议削弱拉萨秩序，并为卡尔藏嘉措承认与清军入藏提供背景影响下一段安排。原有置信注记：被押离拉萨可靠；死亡地点、日期和是否在世的传统说法相互冲突，不能写成确定结论

\paragraph{康熙四十五年（1706年）：第六世达赖喇嘛仓央嘉措被押离拉萨，途中死亡或失踪，其结局与继承合法性成为争议}
边地道路、驻防与多方社群交错的尺度中的\eventanchor{EQ-1706-SIXTH-DALAI-AFTERMATH}{第六世达赖喇嘛仓央嘉措被押离拉萨，途中死亡或失踪，其结局与继承合法性成为争议} 不能离开西藏、和硕特汗国与清来解释。本节点追踪「第六世达赖喇嘛仓央嘉措被押离拉萨，途中死亡或失踪，其结局与继承合法性成为争议」之后可见的余波。拉藏汗以宗教合法性重组拉萨权力，清廷册封语言与西藏各派、蒙古势力认定不一致说明其形成的条件；《清圣祖实录》康熙四十五年西藏条；《清史稿·西藏传》；第巴传记、藏文史书与拉达克、蒙古相关记载提供来源路径，达赖喇嘛地位争议削弱拉萨秩序，并为卡尔藏嘉措承认与清军入藏提供背景则是材料能够支持的近时结果。原有置信注记：被押离拉萨可靠；死亡地点、日期和是否在世的传统说法相互冲突，不能写成确定结论

\paragraph{康熙四十五年（1706年）：第六世达赖喇嘛仓央嘉措被押离拉萨，途中死亡或失踪，其结局与继承合法性成为争议}
\eventanchor{EQ-1706-SIXTH-DALAI-DECISION}{第六世达赖喇嘛仓央嘉措被押离拉萨，途中死亡或失踪，其结局与继承合法性成为争议} 把西藏、和硕特汗国与清放回边地道路、驻防与多方社群交错的尺度的道路、城镇、边界或制度网络。本节点定位围绕「第六世达赖喇嘛仓央嘉措被押离拉萨，途中死亡或失踪，其结局与继承合法性成为争议」形成的处置与调度记录。拉藏汗以宗教合法性重组拉萨权力，清廷册封语言与西藏各派、蒙古势力认定不一致。本段采用《清圣祖实录》康熙四十五年西藏条；《清史稿·西藏传》；第巴传记、藏文史书与拉达克、蒙古相关记载核对其顺序，随后可追到达赖喇嘛地位争议削弱拉萨秩序，并为卡尔藏嘉措承认与清军入藏提供背景。原有置信注记：被押离拉萨可靠；死亡地点、日期和是否在世的传统说法相互冲突，不能写成确定结论

\paragraph{康熙四十五年（1706年）：第六世达赖喇嘛仓央嘉措被押离拉萨，途中死亡或失踪，其结局与继承合法性成为争议}
本节点回看「第六世达赖喇嘛仓央嘉措被押离拉萨，途中死亡或失踪，其结局与继承合法性成为争议」之前已可见的条件，故\eventanchor{EQ-1706-SIXTH-DALAI-PRELUDE}{第六世达赖喇嘛仓央嘉措被押离拉萨，途中死亡或失踪，其结局与继承合法性成为争议} 不只是一个年份标签。西藏、和硕特汗国与清在边地道路、驻防与多方社群交错的尺度承受的压力来自拉藏汗以宗教合法性重组拉萨权力，清廷册封语言与西藏各派、蒙古势力认定不一致。《清圣祖实录》康熙四十五年西藏条；《清史稿·西藏传》；第巴传记、藏文史书与拉达克、蒙古相关记载可回查该节点；达赖喇嘛地位争议削弱拉萨秩序，并为卡尔藏嘉措承认与清军入藏提供背景是其进入后续年序的方式。原有置信注记：被押离拉萨可靠；死亡地点、日期和是否在世的传统说法相互冲突，不能写成确定结论

\paragraph{康熙四十七年（1708年）：卡尔藏嘉措在青海塔尔寺附近被认定为第六世达赖喇嘛转世之一，围绕其身份形成跨区域支持网络}
\eventanchor{EQ-1708-KELZANG}{卡尔藏嘉措在青海塔尔寺附近被认定为第六世达赖喇嘛转世之一，围绕其身份形成跨区域支持网络} 所见的行动者为青海、蒙古诸部与西藏宗教网络，空间尺度为边地道路、驻防与多方社群交错的尺度。本节点叙述该事件本身。第六世达赖喇嘛结局未决和拉藏汗所立人选缺乏广泛承认，青海寺院与蒙古关系提供替代合法性中心使它成为具体事件链的一环；《清圣祖实录》康熙四十七年青海西藏条；《清史稿·西藏传》；塔尔寺藏文档案与第七世达赖喇嘛传与卡尔藏嘉措成为清廷、青海蒙古和西藏各派互动关键人物，宗教承认与军事保护相互牵连。原有置信注记：出生与早期认定大体可靠；具体认定程序、时间和各方承认不可由后世转世谱系倒推

\paragraph{康熙四十七年（1708年）：卡尔藏嘉措在青海塔尔寺附近被认定为第六世达赖喇嘛转世之一，围绕其身份形成跨区域支持网络}
从边地道路、驻防与多方社群交错的尺度看，\eventanchor{EQ-1708-KELZANG-AFTERMATH}{卡尔藏嘉措在青海塔尔寺附近被认定为第六世达赖喇嘛转世之一，围绕其身份形成跨区域支持网络} 留下本节点追踪「卡尔藏嘉措在青海塔尔寺附近被认定为第六世达赖喇嘛转世之一，围绕其身份形成跨区域支持网络」之后可见的余波的材料坐标。青海、蒙古诸部与西藏宗教网络所面对的条件是第六世达赖喇嘛结局未决和拉藏汗所立人选缺乏广泛承认，青海寺院与蒙古关系提供替代合法性中心。《清圣祖实录》康熙四十七年青海西藏条；《清史稿·西藏传》；塔尔寺藏文档案与第七世达赖喇嘛传支持的结果为卡尔藏嘉措成为清廷、青海蒙古和西藏各派互动关键人物，宗教承认与军事保护相互牵连。原有置信注记：出生与早期认定大体可靠；具体认定程序、时间和各方承认不可由后世转世谱系倒推

\paragraph{康熙四十七年（1708年）：卡尔藏嘉措在青海塔尔寺附近被认定为第六世达赖喇嘛转世之一，围绕其身份形成跨区域支持网络}
\eventanchor{EQ-1708-KELZANG-DECISION}{卡尔藏嘉措在青海塔尔寺附近被认定为第六世达赖喇嘛转世之一，围绕其身份形成跨区域支持网络} 记录青海、蒙古诸部与西藏宗教网络在边地道路、驻防与多方社群交错的尺度的一次变化。本节点定位围绕「卡尔藏嘉措在青海塔尔寺附近被认定为第六世达赖喇嘛转世之一，围绕其身份形成跨区域支持网络」形成的处置与调度记录。其发生与第六世达赖喇嘛结局未决和拉藏汗所立人选缺乏广泛承认，青海寺院与蒙古关系提供替代合法性中心相连；《清圣祖实录》康熙四十七年青海西藏条；《清史稿·西藏传》；塔尔寺藏文档案与第七世达赖喇嘛传把事件系回来源，卡尔藏嘉措成为清廷、青海蒙古和西藏各派互动关键人物，宗教承认与军事保护相互牵连显示压力如何向后转移。原有置信注记：出生与早期认定大体可靠；具体认定程序、时间和各方承认不可由后世转世谱系倒推

\paragraph{康熙四十七年（1708年）：卡尔藏嘉措在青海塔尔寺附近被认定为第六世达赖喇嘛转世之一，围绕其身份形成跨区域支持网络}
\eventanchor{EQ-1708-KELZANG-PRELUDE}{卡尔藏嘉措在青海塔尔寺附近被认定为第六世达赖喇嘛转世之一，围绕其身份形成跨区域支持网络} 的地点与行动范围属于边地道路、驻防与多方社群交错的尺度，行动者是青海、蒙古诸部与西藏宗教网络。本节点回看「卡尔藏嘉措在青海塔尔寺附近被认定为第六世达赖喇嘛转世之一，围绕其身份形成跨区域支持网络」之前已可见的条件。第六世达赖喇嘛结局未决和拉藏汗所立人选缺乏广泛承认，青海寺院与蒙古关系提供替代合法性中心。读《清圣祖实录》康熙四十七年青海西藏条；《清史稿·西藏传》；塔尔寺藏文档案与第七世达赖喇嘛传时可见其记录边界，最小后果只能写作卡尔藏嘉措成为清廷、青海蒙古和西藏各派互动关键人物，宗教承认与军事保护相互牵连。原有置信注记：出生与早期认定大体可靠；具体认定程序、时间和各方承认不可由后世转世谱系倒推

\paragraph{康熙五十六年（1717年）：准噶尔军进入拉萨并杀拉藏汗，拉萨政局再次易手}
\eventanchor{EQ-1717-DZUNGAR-LHASA}{准噶尔军进入拉萨并杀拉藏汗，拉萨政局再次易手} 是准噶尔、和硕特汗国与西藏在边地道路、驻防与多方社群交错的尺度留下的编年标记。本节点叙述该事件本身。它从拉藏汗政权缺乏广泛稳定支持，准噶尔借宗教合法性与军事力量越过高原介入走来；《清圣祖实录》康熙五十六年西藏军报；《清史稿·西藏传》；藏文《西藏王臣记》续编与准噶尔相关材料固定可证部分，清廷面临保护卡尔藏嘉措、维护边疆通道和回应准噶尔扩张的复合压力，入藏决策加速构成下一步的可见结果。原有置信注记：入拉萨和拉藏汗死亡可靠；暴力规模、民众态度与准噶尔治理需并读藏文、满文和旅行记录

\paragraph{康熙五十六年（1717年）：准噶尔军进入拉萨并杀拉藏汗，拉萨政局再次易手}
本节点追踪「准噶尔军进入拉萨并杀拉藏汗，拉萨政局再次易手」之后可见的余波：\eventanchor{EQ-1717-DZUNGAR-LHASA-AFTERMATH}{准噶尔军进入拉萨并杀拉藏汗，拉萨政局再次易手}。准噶尔、和硕特汗国与西藏在边地道路、驻防与多方社群交错的尺度的处境可由拉藏汗政权缺乏广泛稳定支持，准噶尔借宗教合法性与军事力量越过高原介入说明。《清圣祖实录》康熙五十六年西藏军报；《清史稿·西藏传》；藏文《西藏王臣记》续编与准噶尔相关材料是本段材料入口，因而只能据此追踪到清廷面临保护卡尔藏嘉措、维护边疆通道和回应准噶尔扩张的复合压力，入藏决策加速。原有置信注记：入拉萨和拉藏汗死亡可靠；暴力规模、民众态度与准噶尔治理需并读藏文、满文和旅行记录

\paragraph{康熙五十六年（1717年）：准噶尔军进入拉萨并杀拉藏汗，拉萨政局再次易手}
\eventanchor{EQ-1717-DZUNGAR-LHASA-DECISION}{准噶尔军进入拉萨并杀拉藏汗，拉萨政局再次易手} 发生在边地道路、驻防与多方社群交错的尺度，牵动准噶尔、和硕特汗国与西藏。本节点定位围绕「准噶尔军进入拉萨并杀拉藏汗，拉萨政局再次易手」形成的处置与调度记录。拉藏汗政权缺乏广泛稳定支持，准噶尔借宗教合法性与军事力量越过高原介入是其结构性条件；《清圣祖实录》康熙五十六年西藏军报；《清史稿·西藏传》；藏文《西藏王臣记》续编与准噶尔相关材料留下的证据把读者带到清廷面临保护卡尔藏嘉措、维护边疆通道和回应准噶尔扩张的复合压力，入藏决策加速。原有置信注记：入拉萨和拉藏汗死亡可靠；暴力规模、民众态度与准噶尔治理需并读藏文、满文和旅行记录

\paragraph{康熙五十六年（1717年）：准噶尔军进入拉萨并杀拉藏汗，拉萨政局再次易手}
\eventanchor{EQ-1717-DZUNGAR-LHASA-PRELUDE}{准噶尔军进入拉萨并杀拉藏汗，拉萨政局再次易手} 对应本节点回看「准噶尔军进入拉萨并杀拉藏汗，拉萨政局再次易手」之前已可见的条件。在边地道路、驻防与多方社群交错的尺度，准噶尔、和硕特汗国与西藏无法绕开拉藏汗政权缺乏广泛稳定支持，准噶尔借宗教合法性与军事力量越过高原介入。依据《清圣祖实录》康熙五十六年西藏军报；《清史稿·西藏传》；藏文《西藏王臣记》续编与准噶尔相关材料，本段把后续影响限定为清廷面临保护卡尔藏嘉措、维护边疆通道和回应准噶尔扩张的复合压力，入藏决策加速。原有置信注记：入拉萨和拉藏汗死亡可靠；暴力规模、民众态度与准噶尔治理需并读藏文、满文和旅行记录

\paragraph{康熙五十七年（1718年）：清军自四川方向入藏，在喀喇乌苏一带覆败，首批救援行动受挫}
\eventanchor{EQ-1718-QING-TIBET-EXPEDITION}{清军自四川方向入藏，在喀喇乌苏一带覆败，首批救援行动受挫}。本节点叙述该事件本身。清、准噶尔与西藏在边地道路、驻防与多方社群交错的尺度的处境首先受清军低估高原道路、气候、补给与准噶尔防御，跨区域指挥链难以及时校正制约。《清圣祖实录》康熙五十七年川藏军报；《清史稿·西藏传》；《康熙朝满文朱批奏折全译》入藏军需奏折及藏文地方记载留下可核的线索；可辨的后果是失败迫使清廷重组由青海、四川等方向的兵站与联盟，为1720年行动积累经验。原有置信注记：远征失利可靠；路线、伤亡和藏地向导问题主要见清方军报，需与藏文地名和地方传统核对

\paragraph{康熙五十七年（1718年）：清军自四川方向入藏，在喀喇乌苏一带覆败，首批救援行动受挫}
\eventanchor{EQ-1718-QING-TIBET-EXPEDITION-AFTERMATH}{清军自四川方向入藏，在喀喇乌苏一带覆败，首批救援行动受挫} 置于边地道路、驻防与多方社群交错的尺度中阅读，本节点追踪「清军自四川方向入藏，在喀喇乌苏一带覆败，首批救援行动受挫」之后可见的余波。清军低估高原道路、气候、补给与准噶尔防御，跨区域指挥链难以及时校正构成行动者清、准噶尔与西藏必须面对的条件。取证从《清圣祖实录》康熙五十七年川藏军报；《清史稿·西藏传》；《康熙朝满文朱批奏折全译》入藏军需奏折及藏文地方记载出发，因而只能把变化谨慎地连到失败迫使清廷重组由青海、四川等方向的兵站与联盟，为1720年行动积累经验。原有置信注记：远征失利可靠；路线、伤亡和藏地向导问题主要见清方军报，需与藏文地名和地方传统核对

\paragraph{康熙五十七年（1718年）：清军自四川方向入藏，在喀喇乌苏一带覆败，首批救援行动受挫}
\eventanchor{EQ-1718-QING-TIBET-EXPEDITION-DECISION}{清军自四川方向入藏，在喀喇乌苏一带覆败，首批救援行动受挫} 标记本节点定位围绕「清军自四川方向入藏，在喀喇乌苏一带覆败，首批救援行动受挫」形成的处置与调度记录，涉及清、准噶尔与西藏。边地道路、驻防与多方社群交错的尺度不是抽象背景：清军低估高原道路、气候、补给与准噶尔防御，跨区域指挥链难以及时校正。《清圣祖实录》康熙五十七年川藏军报；《清史稿·西藏传》；《康熙朝满文朱批奏折全译》入藏军需奏折及藏文地方记载所见的顺序随后指向失败迫使清廷重组由青海、四川等方向的兵站与联盟，为1720年行动积累经验，但原有置信注记：远征失利可靠；路线、伤亡和藏地向导问题主要见清方军报，需与藏文地名和地方传统核对

\paragraph{康熙五十七年（1718年）：清军自四川方向入藏，在喀喇乌苏一带覆败，首批救援行动受挫}
在边地道路、驻防与多方社群交错的尺度，\eventanchor{EQ-1718-QING-TIBET-EXPEDITION-PRELUDE}{清军自四川方向入藏，在喀喇乌苏一带覆败，首批救援行动受挫} 让清、准噶尔与西藏的一个选择或压力有了可查的坐标。本节点回看「清军自四川方向入藏，在喀喇乌苏一带覆败，首批救援行动受挫」之前已可见的条件。前因可写为清军低估高原道路、气候、补给与准噶尔防御，跨区域指挥链难以及时校正；《清圣祖实录》康熙五十七年川藏军报；《清史稿·西藏传》；《康熙朝满文朱批奏折全译》入藏军需奏折及藏文地方记载限定了记录，后续变化是失败迫使清廷重组由青海、四川等方向的兵站与联盟，为1720年行动积累经验。原有置信注记：远征失利可靠；路线、伤亡和藏地向导问题主要见清方军报，需与藏文地名和地方传统核对

\paragraph{康熙五十九年（1720年）：清军及青海蒙古盟军进入拉萨，驱逐准噶尔力量并迎立卡尔藏嘉措为第七世达赖喇嘛}
\eventanchor{EQ-1720-LHASA-RESTORATION}{清军及青海蒙古盟军进入拉萨，驱逐准噶尔力量并迎立卡尔藏嘉措为第七世达赖喇嘛} 所关联的是清、准噶尔、西藏与青海蒙古在边地道路、驻防与多方社群交错的尺度中的行动。本节点叙述该事件本身。它从1718年失败后清廷改进补给并取得青海蒙古支持，准噶尔在拉萨的统治基础也不稳展开，借《清圣祖实录》康熙五十九年西藏军报；《清史稿·西藏传》；满文军机前身档、藏文第七世达赖传与青海蒙古文书进入可检验的年序，并以清廷在拉萨的军事政治存在增强，但西藏政教制度未被单一中央命令取代，后续治理仍待安排影响下一段安排。原有置信注记：进入拉萨和迎立主干可靠；清军、蒙古盟军与西藏地方行动者角色及伤亡不能由凯旋叙事抹平

\paragraph{康熙五十九年（1720年）：清军及青海蒙古盟军进入拉萨，驱逐准噶尔力量并迎立卡尔藏嘉措为第七世达赖喇嘛}
边地道路、驻防与多方社群交错的尺度中的\eventanchor{EQ-1720-LHASA-RESTORATION-AFTERMATH}{清军及青海蒙古盟军进入拉萨，驱逐准噶尔力量并迎立卡尔藏嘉措为第七世达赖喇嘛} 不能离开清、准噶尔、西藏与青海蒙古来解释。本节点追踪「清军及青海蒙古盟军进入拉萨，驱逐准噶尔力量并迎立卡尔藏嘉措为第七世达赖喇嘛」之后可见的余波。1718年失败后清廷改进补给并取得青海蒙古支持，准噶尔在拉萨的统治基础也不稳说明其形成的条件；《清圣祖实录》康熙五十九年西藏军报；《清史稿·西藏传》；满文军机前身档、藏文第七世达赖传与青海蒙古文书提供来源路径，清廷在拉萨的军事政治存在增强，但西藏政教制度未被单一中央命令取代，后续治理仍待安排则是材料能够支持的近时结果。原有置信注记：进入拉萨和迎立主干可靠；清军、蒙古盟军与西藏地方行动者角色及伤亡不能由凯旋叙事抹平

\paragraph{康熙五十九年（1720年）：清军及青海蒙古盟军进入拉萨，驱逐准噶尔力量并迎立卡尔藏嘉措为第七世达赖喇嘛}
\eventanchor{EQ-1720-LHASA-RESTORATION-DECISION}{清军及青海蒙古盟军进入拉萨，驱逐准噶尔力量并迎立卡尔藏嘉措为第七世达赖喇嘛} 把清、准噶尔、西藏与青海蒙古放回边地道路、驻防与多方社群交错的尺度的道路、城镇、边界或制度网络。本节点定位围绕「清军及青海蒙古盟军进入拉萨，驱逐准噶尔力量并迎立卡尔藏嘉措为第七世达赖喇嘛」形成的处置与调度记录。1718年失败后清廷改进补给并取得青海蒙古支持，准噶尔在拉萨的统治基础也不稳。本段采用《清圣祖实录》康熙五十九年西藏军报；《清史稿·西藏传》；满文军机前身档、藏文第七世达赖传与青海蒙古文书核对其顺序，随后可追到清廷在拉萨的军事政治存在增强，但西藏政教制度未被单一中央命令取代，后续治理仍待安排。原有置信注记：进入拉萨和迎立主干可靠；清军、蒙古盟军与西藏地方行动者角色及伤亡不能由凯旋叙事抹平

\paragraph{康熙五十九年（1720年）：清军及青海蒙古盟军进入拉萨，驱逐准噶尔力量并迎立卡尔藏嘉措为第七世达赖喇嘛}
本节点回看「清军及青海蒙古盟军进入拉萨，驱逐准噶尔力量并迎立卡尔藏嘉措为第七世达赖喇嘛」之前已可见的条件，故\eventanchor{EQ-1720-LHASA-RESTORATION-PRELUDE}{清军及青海蒙古盟军进入拉萨，驱逐准噶尔力量并迎立卡尔藏嘉措为第七世达赖喇嘛} 不只是一个年份标签。清、准噶尔、西藏与青海蒙古在边地道路、驻防与多方社群交错的尺度承受的压力来自1718年失败后清廷改进补给并取得青海蒙古支持，准噶尔在拉萨的统治基础也不稳。《清圣祖实录》康熙五十九年西藏军报；《清史稿·西藏传》；满文军机前身档、藏文第七世达赖传与青海蒙古文书可回查该节点；清廷在拉萨的军事政治存在增强，但西藏政教制度未被单一中央命令取代，后续治理仍待安排是其进入后续年序的方式。原有置信注记：进入拉萨和迎立主干可靠；清军、蒙古盟军与西藏地方行动者角色及伤亡不能由凯旋叙事抹平